\scanpage{31}
% archive foliation: 30
$J_0$ de $J$
\[
Y'_{J_0}=\coprod_{j\in J_0}\!{}_Y\,\unsure{Y'_{J_0}|Y'_j},
\]
et $Y'_J=\varinjlim Y'_{J_0}$, on aura $Y\longrightarrow Y'\in\underline W$. Donc il y aura quand $=$ surjectivité! Mais soient maintenant
\diagram{\begin{tikzcd}[column sep=large]
f,g:X\arrow[r,shift left=.5ex,"\unsure{f_0|f}"]\arrow[r,shift right=.5ex,"g"']&\deleted{Y'}\inserted{Y}
\end{tikzcd}}
qui \deleted{\unsure{deviennent}} sont tels que $\unsure{\sigma|\gamma}(f_j)=\unsure{\sigma|\gamma}(g_j)$ $\forall j\in J$, a-t-on $\unsure{\sigma|\gamma}(f)=\unsure{\sigma|\gamma}(g)$? Mais oui, par le $=$ argument --- condensé dans un

\margin{\uuline{bof!}}

\uline{Lemme} Soit $Y\xrightarrow{\alpha_j}Y'_j$ \deleted{de} une \unsure{foule|famille} d'équivs. faibles, alors il existe $Y\xrightarrow{\alpha}Y'$ une équiv. faible \deleted{telle que} \inserted{$Y\lhook\joinrel\longrightarrow Y'$} (imm. ouverte ou fermée si on y tient) \deleted{telle que} des $\beta_j:Y'_j\longrightarrow Y'$ tels que $\beta_j\alpha_j=\alpha$.

\margin{\diagram{\begin{tikzcd}[column sep=large,row sep=large]
Y\arrow[r,hook,"\bar\alpha_j"]\arrow[dr,"\alpha"']&\overline{Y'_j}\arrow[r,shift left=.5ex,"i_j"]\arrow[d]&Y'_j\arrow[l,dashed,shift left=.5ex,"r_j"]\arrow[dl,dashed,"\beta_j"]\\
&Y'&
\arrow[from=1-1,to=2-2,phantom,"\mathrm{comm}"{pos=.55}]
\arrow[from=1-3,to=2-2,phantom,"\mathrm{comm}"'{pos=.6}]
\end{tikzcd}}
$r_ji_j=\mathrm{id}_{\overline{Y'_j}}$}

(De plus, si les $\alpha_j$ sont des imm. ouvertes, resp. fermées, on peut prendre pour $Y'$ la somme amalgamée des $Y'_j$.)

On est ramené à prouver ceci: Soient $f,g:X=\coprod X_j\longrightarrow Y$, tels que $f_j\sim g_j$ $\forall j\in J$, alors $\unsure{\sigma|\gamma}(f)=\unsure{\sigma|\gamma}(g)$. (Mais attention on n'aura pas $f\sim g$!) \uline{Il faut une condition du type ``\unsure{Kan}'' sur $Y$}, ou plutôt un $Y'$, où $(Y\longrightarrow Y')\in\underline W$.
